\documentclass[12pt,a4paper]{article}
\usepackage[T2A]{fontenc}
\usepackage[utf8]{inputenc}
\usepackage[russian]{babel}
\usepackage{geometry}
\usepackage{booktabs}
\usepackage{hyperref}
\usepackage{amsmath}
\geometry{margin=2cm}

\title{SET 9. Задача A1: анализ строковых сортировок}
\author{Логинов Николай}
\date{Май 2026}

\begin{document}
\maketitle

\section*{Цель работы}
Цель эксперимента --- сравнить стандартные алгоритмы сортировки строк с алгоритмами,
которые используют посимвольную структуру строк: тернарный String Quick Sort,
String Merge Sort с учетом LCP, MSD Radix Sort и MSD Radix Sort с переключением на
String Quick Sort при размере подмассива меньше 74.

\section*{Подготовка данных}
Для генерации данных используется класс \texttt{StringGenerator}. Алфавит содержит
74 символа: латинские буквы верхнего и нижнего регистра, цифры и символы
\texttt{!@\#\%:;\^{}\&*()-.}. Длины строк выбираются случайно от 10 до 200 символов.
Для каждого типа данных сначала генерируется массив из 3000 строк, затем из него
берутся префиксы размеров 100, 200, \ldots, 3000.

Рассматривались четыре типа массивов:
\begin{itemize}
  \item случайные массивы;
  \item обратно отсортированные массивы;
  \item почти отсортированные массивы, полученные небольшим числом перестановок;
  \item массивы с искусственно добавленным общим префиксом.
\end{itemize}

\section*{Методика измерений}
Измерения выполняются классом \texttt{StringSortTester}. Для каждого запуска фиксируются
среднее время работы в микросекундах и количество посимвольных операций. Для стандартных
Quick Sort и Merge Sort считаются сравнения символов внутри лексикографического сравнения.
Для MSD Radix Sort дополнительно естественно учитывать обращения к очередному символу,
поскольку алгоритм не сравнивает пары строк напрямую, а классифицирует строки по символам.

Команда для воспроизведения полного эксперимента:
\begin{verbatim}
g++ -std=c++17 -O2 A1_experiment.cpp -o A1_experiment
./A1_experiment > results.csv
\end{verbatim}

Для быстрой проверки корректности:
\begin{verbatim}
./A1_experiment --quick > quick_results.csv
\end{verbatim}

\section*{Теоретические ожидания}
Для обычных сортировок сравнением число сравнений строк имеет порядок
$O(n \log n)$, но каждое сравнение может требовать просмотра общего префикса,
поэтому фактическая стоимость зависит от суммарной длины просмотренных префиксов.
На массивах с длинными совпадающими префиксами это особенно заметно.

Тернарный String Quick Sort группирует строки по символу на текущей позиции и
переходит к следующему символу только для равной группы. В среднем он близок к
$O(W + n \log n)$, где $W$ --- суммарное число просмотренных символов.

MSD Radix Sort распределяет строки по символу текущей позиции и рекурсивно сортирует
только непустые корзины. Его ожидаемая стоимость близка к $O(W + Rn)$ на верхних
уровнях, где $R$ --- размер алфавита. Переключение на String Quick Sort при малых
подмассивах уменьшает накладные расходы на массивы счетчиков.

\section*{Выводы по эксперименту}
После запуска эксперимента таблица \texttt{results.csv} содержит строки вида:
\begin{verbatim}
dataset,size,algorithm,time_us,char_ops,sorted
\end{verbatim}

Ожидаемый качественный результат:
\begin{itemize}
  \item стандартные Quick Sort и Merge Sort проигрывают на данных с длинными общими
        префиксами, так как многократно сравнивают одни и те же начальные символы;
  \item String Quick Sort обычно уменьшает число повторных посимвольных сравнений;
  \item MSD Radix Sort хорошо работает на случайных строках, но имеет заметные
        накладные расходы на маленьких подмассивах;
  \item гибрид MSD Radix + String Quick Sort обычно быстрее чистого MSD Radix на
        небольших корзинах.
\end{itemize}

\section*{Материалы для сдачи}
\begin{itemize}
  \item Реализация инфраструктуры: \texttt{A1\_experiment.cpp}.
  \item ID посылки A1m: 374545278.
  \item ID посылки A1q: 374545410.
  \item ID посылки A1r: 374545486.
  \item ID посылки A1rq: 374545543.
  \item Ссылка на публичный репозиторий с результатами: \url{https://github.com/elegant784seat/set9.git}.
\end{itemize}

\end{document}
